\section{Limitations}
\label{sec:limitations}

The results are deliberately narrow. This section lists what the study cannot support, so that
the claims in \Cref{tab:claims} are not read more broadly than the evidence allows.

\paragraph{Statistical power.}
Most comparisons use two or three seeds, and we do not compute significance tests. The
seed-to-seed spread on the aggregate task is larger than the effect being measured, which is
why \Cref{sec:results-gap-bc} calls it inconclusive rather than positive. A single seed of the
largest Gap~D arm would not have been convincing; the three-seed presence result is credible
because every seed is far above chance and the arms are otherwise identical. The Gap~A result
is a controlled ablation with frozen parameters, so seed noise is small, but it remains a
single task.

\paragraph{Construct validity of the synthetic tasks.}
MQAR, delayed recall, aggregate count, and presence recall each isolate one operator. That is
their purpose, but isolation is also their weakness: a mechanism that works on a synthetic
binding task need not be the mechanism a language model uses. The presence task in particular
uses frozen embeddings and untied heads to remove a shortcut; that control changes the
optimisation problem relative to a normal language model, so the $0.977$ accuracy should be
read as evidence about the read operator, not as a prediction of LM behaviour.

\paragraph{No matched transformer language-model baseline.}
This is the most important omission. The transformer in the experiment ladder is a width-matched
control for synthetic recall and was never trained on TinyStories. It therefore cannot answer
whether SigmaStorm is better or worse than a transformer at language modeling. Such a comparison
requires the same tokenizer, corpus, token budget, optimiser, evaluation split, and roughly
matched parameters or FLOPs. Until it exists, no capability comparison with transformers is
supported, including comparisons implied by the architecture's motivation.

\paragraph{Scale.}
The largest model is 21M parameters, trained for roughly 348M tokens on TinyStories. There are
no scaling curves, no runs at 100M or 1B parameters, and no data-mix study. TinyStories is a
restricted corpus with simple syntax; success there says little about natural text. The
tokenizer is an 8K byte-level BPE, so the reported perplexity is not comparable to perplexities
computed with larger subword vocabularies.

\paragraph{Long context.}
Every operator is chunk-invariant and the state is constant in $T$, but the largest trained run
used 1{,}024-token windows with a differentiable horizon of $4\times256$ tokens in the reported
configuration. The 256K design target has not been trained or evaluated. Specifically unmeasured:
retrieval over tens of thousands of tokens, stability of the energy-controlled latent over long
streams, the interaction between the summary ring and the delta state at scale, and the quality
of generation after chunked prefill. The chunk-invariance proof is a statement about
implementation correctness, not about learned long-context ability.

\paragraph{Efficiency measurements.}
The cost table reports CPU forward times with two threads and batch size one. Those numbers
measure dense PyTorch implementations without fused kernels for the chunkwise delta solve, the
slot consolidation, or the summary scan. They are useful as a relative sanity check and as
evidence that the operators are linear in $T$; they are not deployment throughput and they do
not constitute a fair efficiency comparison against optimized attention kernels.

\paragraph{Hardware and reruns.}
The 21M language-model run used two 16\,GB T4 GPUs. The mechanism and Gap~D experiments were
re-run on CPU after the weekly GPU quota was exhausted, with identical data, seeds, optimiser,
and step counts. This is honest but not ideal: CPU and GPU floating-point behaviour differ, and
the re-runs do not exercise the distributed training path. The absence of a multi-GPU scaling
result is a gap in the engineering evidence, not just in the science.

\paragraph{Evaluation breadth.}
There are no downstream benchmarks, no zero-shot evaluations, no instruction tuning, and no
human evaluation. The 21M model is a prototype that writes on-genre prose; it is not evaluated
as an assistant. Claims about usefulness, safety, or alignment are outside the scope of this
report and are not made.

\paragraph{Known negative and unstable results.}
State expansion reduced accuracy at 66 bindings, and the summary benefit changed sign between
runs. Both are reported as results rather than suppressed. They limit the architectural story:
the demonstrated recall mechanisms are the delta rule, the decay spectrum, and the gated slot
read, not the additional state matrices or the summary ring. A reader should not infer that
every component of the block contributes measured value.

\paragraph{Reproducibility caveats.}
All headline numbers are produced by scripts in the repository, and the raw artifacts are
checked in. Reproducing them still requires the same package versions and, for the 21M run, the
same two-GPU configuration. We report seeds and step counts, but we do not claim bitwise
determinism across hardware. The paper's evidence policy is traceability to an artifact, not
independent replication.
